\section{Worked Examples and Probe Evidence}

\tool is evaluated as an interface resource: can one adapter contract turn
heterogeneous \sid exports into comparable artifact profiles? GRID supplies a
released GRID/RQ-KMeans-style module; ReSID supplies processed Musical artifacts
and a bounded GAOQ export path. Sanity rows and controlled mechanism probes then
calibrate whether diagnostic differences are structural, behavioral, or
capacity-driven.

\paragraph{Worked-example protocol.}
The case study uses the same 23,742 Musical items: a controlled GRID/RQ-KMeans
feature-text export from ReSID text features, plus a bounded ReSID
FAMAE-to-GAOQ export. This supports same-item artifact comparison, not faithful
reproduction of TIGER, GRID, or full ReSID training.

\begin{table}[t]
\caption{SID diagnostics on Amazon Reviews 2023 Musical Instruments. Rows
follow the diagnostic path: RQ-style variants, ReSID contrast, then controls.
Bold rows are named-tokenizer export paths; italics are controls.}
\label{tab:musical-diagnostic}
\scriptsize
\setlength{\tabcolsep}{1.3pt}
\begin{tabular}{@{}>{\raggedright\arraybackslash}p{0.19\columnwidth}>{\raggedleft\arraybackslash}p{0.15\columnwidth}>{\raggedleft\arraybackslash}p{0.10\columnwidth}>{\raggedleft\arraybackslash}p{0.10\columnwidth}>{\raggedleft\arraybackslash}p{0.10\columnwidth}>{\raggedleft\arraybackslash}p{0.24\columnwidth}@{}}
\toprule
Artifact & Unique SIDs & D2 & D3 & D4 & D5 \\
\midrule
\textbf{GRID ft} & 3,749 & 0.977 & 0.055 & 0.370 & 64/3.4k/3.7k \\
GRID ft-cap & 9,874 & 0.779 & 0.080 & 0.639 & 32/9300/9874 \\
RQ-min ref & 17,247 & 0.440 & 0.065 & 0.883 & 32/2368/17.2k \\
\textbf{ReSID}\textsuperscript{\dag} & 23,742 & 0.000 & 0.154 & 1.000 & 32/1280/23.7k \\
\emph{Cat-prefix}\textsuperscript{\dag} & 23,742 & 0.000 & 0.447 & 1.000 & 30/83/313/23.7k \\
\emph{Pop-balanced} & 22,707 & 0.086 & 0.303 & 0.962 & 4/1024/22.7k \\
\emph{Hash-collide} & 256 & 1.000 & 0.004 & 0.032 & 256/256/256/256 \\
\bottomrule
\end{tabular}
\end{table}

\begin{table}[t]
\caption{Controlled mechanism probes. Each compact contrast activates a known
artifact pressure and checks whether the targeted diagnostic separates it.}
\label{tab:mechanism-probes}
\scriptsize
\setlength{\tabcolsep}{2pt}
\begin{tabularx}{\columnwidth}{@{}>{\raggedright\arraybackslash}p{0.29\columnwidth}>{\centering\arraybackslash}p{0.11\columnwidth}>{\raggedright\arraybackslash}X>{\raggedright\arraybackslash}X@{}}
\toprule
Probe & D & Without probe & With probe \\
\midrule
Qualified aliasing & D2 & hash 1.19x & co-occur 3.86x \\
Capacity budget & D1/D4 & head 1.000 & tail 0.028 \\
Variable depth & D5 & max-depth 12,010 & active 7,914 \\
\bottomrule
\end{tabularx}
\end{table}

Table~\ref{tab:musical-diagnostic} shows how two exports become comparable
artifact-profile rows under one schema. Its order is intentional: the
RQ-style rows are adjacent so that capacity and aliasing changes can be read
before the ReSID contrast and the sanity controls.
The GRID ft row is intentionally labeled as feature-text: it uses processed
feature text from the ReSID artifact path and the released residual
k-means module, not a faithful raw-text TIGER reproduction. Under that
controlled setup, the row exposes high aliasing pressure and limited tail
capacity. The ReSID/GAOQ row exports one full code per item in this bounded
run, yielding zero full-code collisions and higher D3-L1 collaborative prefix
recovery. The daggered rows make the structural floor explicit: collision-free
full codes are not by themselves evidence of better recommendation quality.
The capacity-matched GRID row reopens the strongest alternative explanation:
when the same feature-text path uses a larger
32/1280/1280 budget, unique full codes rise to 9,874 and D2 aliasing drops to
0.779, but the artifact still does not approach the structurally item-unique
ReSID or category-prefix rows. Thus capacity explains part of the original GRID
pressure. The ablation is not a full item-unique leaf match; additional
capacity could reduce aliasing further, so the contrast remains an artifact
profile rather than a method ranking. The RQ-min row is included for a
different reason: it is a minimal residual-quantization reference adapter that
passes the same validator on the full Musical item universe. Its role is to
show that \tool accepts an independent SID-generation code path; it is not
third named-method coverage and it does not reproduce TIGER, GRID, or ReSID.

\paragraph{Diagnostic findings.}
The worked example supports one central finding and two supporting ones. The
central finding is that addressability is not behavioral prefix alignment. The
ReSID row is structurally item-unique; the category-prefix control is also
alias-free because it uses metadata categories as early levels and a
within-prefix item index as the leaf. It is a taxonomy control rather than a
learned tokenizer. D3 exposes the surprising part of the profile: the
non-learned category-prefix row recovers
more level-1 co-occurrence neighbors (0.447) than ReSID (0.154) or either
GRID row (0.055--0.080). The right interpretation is not that category IDs
are better tokenizers. It is that learned or item-unique addresses do not
guarantee behaviorally aligned prefixes. This failure mode is visible only when
the mapping is inspected directly.

The first supporting finding is intervention-style: a capacity change fixes one
failure mode without repairing another. Increasing the GRID feature-text budget to
32/1280/1280 raises unique full codes from the 3.7--4.0k range to 9,874 and
reduces D2 aliasing from 0.977 to 0.779. That ablation weakens a simple
``too little capacity'' objection, but it does not make the row item-unique.
The RQ-style rows expose a sharper diagnostic dissociation: D2 aliasing falls
from 0.977 (GRID ft) to 0.779 (GRID ft-cap) and 0.440 (RQ-min), while D3
stays in the narrow 0.055--0.080 range. Capacity and aliasing can improve
without producing behaviorally aligned prefixes, which is why these probes are
reported separately. Table~\ref{tab:mechanism-probes} gives the same point as
controlled mechanism checks: qualified aliasing separates co-occurrence aliases
from hash aliases, capacity probes separate head and tail allocation, and
variable-depth probes separate maximum-depth from realized prefix cost.

We also test whether the central D3 finding is fragile. It is not confined to
the Musical table: on an All\_Beauty 20k vertical, a
coarse category-prefix control reaches D3-L1 0.968 while GRID/RQ-KMeans
feature-text rows are 0.081, 0.087, and 0.090 across seeds 42--44. The
coarse-category flag is essential--this is diagnostic portability evidence, not
a claim that taxonomy IDs are better tokenizers. The unusually high
All\_Beauty value is itself part of the diagnostic signal: D3 distinguishes a
vertical where the available coarse taxonomy is strongly aligned with observed
co-occurrence from Musical, where the same category-prefix control is only
moderately aligned. A plausible explanation is that Beauty co-purchase sessions
often stay within narrow product-type and routine categories, while Musical
accessory and instrument co-purchases can cross coarser taxonomy boundaries.
To keep D3 from being only an internal score, we also report bounded ranking
context: a 5,000-user Musical check uses SID prefixes only as candidate sets
and applies the same train-only co-occurrence/popularity reranker to all six
artifact/control rows; D3 correlates with candidate recall (Spearman 0.943) and
fixed-reranker Recall@20 (0.886). A bounded All\_Beauty temporal-LOO
replication is consistent in direction across four rows: D3 is monotonic with
candidate recall and fixed-reranker Recall@20, but the small row count and
constructed split keep it as supplementary directional evidence. Finally, the
GRID export path is not confined to one vertical: a real Sports 20k GRID export
has complete metadata and interaction joins, D3-L1 0.055, duplicate-\sid rate
0.592, and D5 prefixes 128/7986/8165. Sports proxy rows remain supplement-only
and do not add named method coverage. DACT D6 and MovieLens schema rows
exercise extension and portability paths, but the main paper claim remains the
Musical worked example.
